\documentclass[11pt,a4paper]{article}
\usepackage{layout}
\begin{document}

% Page 1: restore the approved 2026 chapter opener.
\begin{canvas}
\pdfbookmark[0]{第一章 从开放协作到公共价值}{chapter1}
\node[anchor=north west,inner sep=0pt,outer sep=0pt] at (current page.north west){\includegraphics[width=210mm,height=297mm]{.build/assets/chapter1-cover.pdf}};
\Box{27}{227}{161}{\sffamily\bfseries\fontsize{11.5}{16}\selectfont\color{teal}AI 重塑开源，开源塑造 AI}
\end{canvas}

% Page 2: AI -> open collaboration, with distinct external windows.
\begin{canvas}
\Header{1.1 · 双向重塑}{blue}\Footer{02}
\Sec{blue}{1.1　人工智能时代的开源协作：双向重塑}{section11}
\Title{AI 改变创造与协作的节奏}
\Subhead{59}{blue}{从活动增长，进入生产方式与协作责任的变化}
\Body{22}{72}{166}{人工智能将生成、自动化和智能体协作带入开发过程，也扩展了代码之外的模型与数据供给。GitHub 与 Hugging Face（HF）的官方观察呈现出不同侧面的增长。规模变化提出了新的研究问题：谁在创造，谁负责审阅与集成，谁承担持续维护？}
\Label{22}{97}{166}{开发活动：GitHub 的三个年度观察}
\Box{22}{106}{52}{\sffamily\bfseries\fontsize{10}{15}\selectfont\color{blue}代码提交}
\Box{22}{114}{52}{\sffamily\fontsize{7.5}{10}\selectfont\color{muted}单位：十亿次}
\Rule{22}{160}{69}{line}
\Rect{25}{151.2536443148688}{10}{8.746355685131196}{blue!22}
\Box{22}{144.2536443148688}{16}{\sffamily\bfseries\fontsize{8.7}{12}\selectfont\color{navy}\centering 10.0}
\Box{22}{162}{16}{\sffamily\fontsize{7.8}{10}\selectfont\color{muted}\centering 2024}
\Rect{40}{148.71720116618076}{10}{11.282798833819243}{blue!47}
\Box{37}{141.71720116618076}{16}{\sffamily\bfseries\fontsize{8.7}{12}\selectfont\color{navy}\centering 12.9}
\Box{37}{162}{16}{\sffamily\fontsize{7.8}{10}\selectfont\color{muted}\centering 2025}
\Rect{55}{130.0}{10}{30.000000000000004}{blue}
\Box{52}{123.0}{16}{\sffamily\bfseries\fontsize{8.7}{12}\selectfont\color{navy}\centering 34.3}
\Box{52}{162}{16}{\sffamily\fontsize{7.8}{10}\selectfont\color{muted}\centering 2026}
\Box{79}{106}{52}{\sffamily\bfseries\fontsize{10}{15}\selectfont\color{blue}Actions 作业}
\Box{79}{114}{52}{\sffamily\fontsize{7.5}{10}\selectfont\color{muted}单位：十亿次}
\Rule{79}{160}{126}{line}
\Rect{82}{154.7136563876652}{10}{5.286343612334802}{blue!22}
\Box{79}{147.7136563876652}{16}{\sffamily\bfseries\fontsize{8.7}{12}\selectfont\color{navy}\centering 8.0}
\Box{79}{162}{16}{\sffamily\fontsize{7.8}{10}\selectfont\color{muted}\centering 2024}
\Rect{97}{150.7488986784141}{10}{9.251101321585903}{blue!47}
\Box{94}{143.7488986784141}{16}{\sffamily\bfseries\fontsize{8.7}{12}\selectfont\color{navy}\centering 14.0}
\Box{94}{162}{16}{\sffamily\fontsize{7.8}{10}\selectfont\color{muted}\centering 2025}
\Rect{112}{130.0}{10}{30.0}{blue}
\Box{109}{123.0}{16}{\sffamily\bfseries\fontsize{8.7}{12}\selectfont\color{navy}\centering 45.4}
\Box{109}{162}{16}{\sffamily\fontsize{7.8}{10}\selectfont\color{muted}\centering 2026}
\Box{136}{106}{52}{\sffamily\bfseries\fontsize{10}{15}\selectfont\color{blue}合并 PR}
\Box{136}{114}{52}{\sffamily\fontsize{7.5}{10}\selectfont\color{muted}单位：十亿次}
\Rule{136}{160}{183}{line}
\Rect{139}{150.72307692307692}{10}{9.276923076923078}{blue!22}
\Box{136}{143.72307692307692}{16}{\sffamily\bfseries\fontsize{8.7}{12}\selectfont\color{navy}\centering 0.402}
\Box{136}{162}{16}{\sffamily\fontsize{7.8}{10}\selectfont\color{muted}\centering 2024}
\Rect{154}{147.58461538461538}{10}{12.415384615384616}{blue!47}
\Box{151}{140.58461538461538}{16}{\sffamily\bfseries\fontsize{8.7}{12}\selectfont\color{navy}\centering 0.538}
\Box{151}{162}{16}{\sffamily\fontsize{7.8}{10}\selectfont\color{muted}\centering 2025}
\Rect{169}{130.0}{10}{30.0}{blue}
\Box{166}{123.0}{16}{\sffamily\bfseries\fontsize{8.7}{12}\selectfont\color{navy}\centering 1.30}
\Box{166}{162}{16}{\sffamily\fontsize{7.8}{10}\selectfont\color{muted}\centering 2026}
\Box{22}{183}{166}{\sffamily\bfseries\fontsize{10.5}{15}\selectfont\color{teal}HF 开放资产：模型、数据与应用空间同步扩张}
\Box{22}{195}{52}{\sffamily\bfseries\fontsize{9.5}{14}\selectfont\color{navy}公共模型仓库}
\Box{54}{195}{19}{\sffamily\bfseries\fontsize{9.5}{14}\selectfont\color{teal}\raggedleft +22\%}
\Rect{22}{208}{29.554054054054053}{4}{teal!25}
\Box{53.55405405405405}{206.5}{14}{\sffamily\fontsize{8.5}{11}\selectfont\color{navy}243}
\Rect{22}{218}{36.0}{4}{teal}
\Box{60.0}{216.5}{14}{\sffamily\fontsize{8.5}{11}\selectfont\color{navy}296}
\Box{79}{195}{52}{\sffamily\bfseries\fontsize{9.5}{14}\selectfont\color{navy}数据集}
\Box{111}{195}{19}{\sffamily\bfseries\fontsize{9.5}{14}\selectfont\color{teal}\raggedleft +41\%}
\Rect{79}{208}{25.596}{4}{teal!25}
\Box{106.596}{206.5}{14}{\sffamily\fontsize{8.5}{11}\selectfont\color{navy}71.1}
\Rect{79}{218}{36.0}{4}{teal}
\Box{117.0}{216.5}{14}{\sffamily\fontsize{8.5}{11}\selectfont\color{navy}100}
\Box{136}{195}{52}{\sffamily\bfseries\fontsize{9.5}{14}\selectfont\color{navy}Spaces}
\Box{168}{195}{19}{\sffamily\bfseries\fontsize{9.5}{14}\selectfont\color{teal}\raggedleft +44\%}
\Rect{136}{208}{25.0}{4}{teal!25}
\Box{163.0}{206.5}{14}{\sffamily\fontsize{8.5}{11}\selectfont\color{navy}100}
\Rect{136}{218}{36.0}{4}{teal}
\Box{174.0}{216.5}{14}{\sffamily\fontsize{8.5}{11}\selectfont\color{navy}144}
\Box{22}{228}{166}{\sffamily\fontsize{7.8}{11}\selectfont\color{muted}浅色为来源所列期初，深色为期末；单位：万个。增幅按展示数值计算。}

\Body{22}{239}{166}{GitHub 还观察到，2026 年 9 月逾 40\% 的 PR 由开发者借助智能体发起。生成变快之后，审阅、测试、集成与维护如何跟进，成为理解开源协作的重要问题。}
\Note{22}{260}{166}{图 1-1　来源：\href{\GH}{GitHub Octoverse 2026}、\href{\HF}{HF Summer 2026}。GitHub 各年为截至 9 月 30 日的 12 个月，含公开与私有仓库；HF 为夏季报告所列公共资源存量，具体快照日未明示。各面板独立标度；两平台不合并计算。}
\end{canvas}

% Page 3: complementary evidence of reuse, usage and cross-border continuation.
\begin{canvas}
\Header{1.1 · 双向重塑}{blue}\Footer{03}
\Kicker{blue}{1.1 / OPEN COLLABORATION SHAPES AI}
\Title{开源改变 AI 的扩散与演进}
\Subhead{59}{blue}{中国模型与全球社区：开放供给如何形成全球接续}
\Body{22}{72}{166}{Qwen 的广泛派生与 DeepSeek 的大型模型使用，展现了不同的生态位置。开放影响还体现在模型输出、数据与工程组件被其他机构继续使用，形成跨组织、跨地域的再创造。}
\Rect{22}{99}{80}{64}{blue!5}
\Rect{22}{99}{80}{1}{blue}
\Box{26}{103}{72}{\sffamily\bfseries\fontsize{10.3}{15}\selectfont\color{blue}Qwen｜广泛的社区派生}
\Rect{108}{99}{80}{64}{teal!5}
\Rect{108}{99}{80}{1}{teal}
\Box{112}{103}{72}{\sffamily\bfseries\fontsize{10.3}{15}\selectfont\color{teal}DeepSeek｜大型模型与推理使用}
\Box{26}{115}{72}{\sffamily\bfseries\fontsize{24}{28}\selectfont\color{blue}151,448}
\Box{26}{129}{72}{\sffamily\fontsize{9}{13}\selectfont\color{navy}HF 报告记录的社区衍生仓库}
\Box{26}{141}{72}{\sffamily\fontsize{8.4}{12}\selectfont\color{muted}28,531 个 GGUF 转换中，官方发布 54 个；}
\Box{26}{149}{72}{\sffamily\fontsize{8.4}{12}\selectfont\color{muted}约 99.8\% 来自其他账号。}
\Box{112}{116}{34}{\sffamily\bfseries\fontsize{23}{28}\selectfont\color{teal}47\%}
\Box{151}{116}{33}{\sffamily\bfseries\fontsize{23}{28}\selectfont\color{teal}31.1\%}
\Box{112}{129}{34}{\sffamily\fontsize{8.5}{12}\selectfont\color{navy}大模型下载份额}
\Box{151}{129}{33}{\sffamily\fontsize{8.5}{12}\selectfont\color{navy}推理词元份额}
\Box{112}{141}{72}{\sffamily\fontsize{8.4}{12}\selectfont\color{muted}总参数 2500 亿及以上模型；}
\Box{112}{149}{72}{\sffamily\fontsize{8.4}{12}\selectfont\color{muted}OpenRouter 月度前十开放模型。}
\Box{22}{169}{166}{\sffamily\bfseries\fontsize{10.5}{15}\selectfont\color{navy}从模型输出，到开放数据，再到后续模型}
\Rect{22}{181}{43}{36}{blue!6}
\Rect{75}{181}{64}{36}{teal!6}
\Rect{149}{181}{39}{36}{gold!6}
\Box{26}{185}{35}{\sffamily\bfseries\fontsize{11}{15}\selectfont\color{blue}DeepSeek-R1}
\Box{26}{195}{35}{\sffamily\fontsize{8.5}{12}\selectfont\color{muted}生成推理等样本}
\Box{79}{184}{56}{\sffamily\bfseries\fontsize{9.5}{14}\selectfont\color{teal}NVIDIA 开放后训练数据}
\Box{79}{194}{56}{\sffamily\bfseries\fontsize{18}{22}\selectfont\color{teal}121.3 万条}
\Box{79}{207}{56}{\sffamily\fontsize{7.8}{11}\selectfont\color{muted}明确列为 R1 生成的样本}
\Box{153}{186}{31}{\sffamily\bfseries\fontsize{9.2}{14}\selectfont\color{gold}Llama-Nemotron}
\Box{153}{197}{31}{\sffamily\bfseries\fontsize{9.2}{14}\selectfont\color{gold}系列模型}
\LinkLine{66}{199}{73}{199}{teal}
\LinkLine{140}{199}{147}{199}{gold}
\Box{22}{220}{166}{\sffamily\bfseries\fontsize{8.8}{13}\selectfont\color{navy}图 1-2　派生、使用与数据接续，共同揭示开放影响}
\Box{22}{227}{166}{\sffamily\fontsize{7.8}{11}\selectfont\color{muted}衍生仓库不等于原创模型；合成样本复用不表示权重继承，样本数量不代表性能贡献比例。}
\Box{22}{238}{166}{\sffamily\bfseries\fontsize{9.5}{14}\selectfont\color{teal}开放继续延伸至基础设施}
\Box{22}{245}{166}{\sffamily\fontsize{8.7}{12}\selectfont\color{ink}DeepSeek 开放计算内核、通信、并行训练与存储组件，为模型之外的体系协作提供基础。}

\Note{22}{256}{166}{来源：\href{\HF}{HF Summer 2026}；\href{\ATOM}{ATOM 论文 v1（2026-04）}图 9、12；\href{\NV}{NVIDIA（2025-05-14）}表 2；\href{\DSINFRA}{DeepSeek 官方开放基础设施索引}。}
\Note{22}{265}{166}{口径：Qwen 为 HF 夏季报告观察；47\% 为 ATOM 所追踪 250B+ 模型累计下载份额，按总参数分组；31.1\% 为论文标注“2026 年 1 月”观察点的 OpenRouter 每月前十开放模型词元份额。非全球市场份额，各指标不直接比较；NVIDIA 案例属历史背景，均不并入本报告主期总量。}
\end{canvas}

% Page 4: approved X1-X4 / D01-D10 identity.
\begin{canvas}
\Header{1.2 · 发展目标}{teal}\Footer{04}
\Sec{teal}{1.2　开源发展目标：从开放资源到公共能力}{section12}
\Title{四项发展目标，组织共同问题}
\Subhead{59}{teal}{以人才、供给、连接与治理观察发展质量}
\Body{22}{72}{166}{AI 与开源的双向变化，使评价需要同时关心参与机会、开放供给、全球接续和长期治理。开源发展目标（OSDGs）以四个互补支柱、十个观察维度组织这些问题，为不同地区与主体提供共同的研究语言。}
\Rule{22}{96}{188}{line}
\Box{22}{99}{15}{\sffamily\bfseries\fontsize{13}{18}\selectfont\color{blue}X1}
\Box{40}{99}{148}{\sffamily\bfseries\fontsize{12}{18}\selectfont\color{blue}人才规模与多样性}
\Box{40}{109}{148}{\sffamily\fontsize{10}{15}\selectfont\color{ink}谁能够进入开放创新，并成长为持续贡献者？}
\definecolor{dD01}{HTML}{155FA8}
\Circle{41.3}{124}{1.1}{dD01}
\Box{45}{120}{44}{\sffamily\fontsize{8.4}{12}\selectfont\color{dD01}D01 规模与密度}
\definecolor{dD02}{HTML}{087C99}
\Circle{90.3}{124}{1.1}{dD02}
\Box{94}{120}{44}{\sffamily\fontsize{8.4}{12}\selectfont\color{dD02}D02 生态多样性}
\Rule{22}{131}{188}{line}
\Box{22}{134}{15}{\sffamily\bfseries\fontsize{13}{18}\selectfont\color{teal}X2}
\Box{40}{134}{148}{\sffamily\bfseries\fontsize{12}{18}\selectfont\color{teal}技术产出与开放供给}
\Box{40}{144}{148}{\sffamily\fontsize{10}{15}\selectfont\color{ink}哪些代码、模型与数据能被获取、理解并继续创造？}
\definecolor{dD03}{HTML}{007D78}
\Circle{41.3}{159}{1.1}{dD03}
\Box{45}{155}{44}{\sffamily\fontsize{8.4}{12}\selectfont\color{dD03}D03 代码产出}
\definecolor{dD04}{HTML}{267D4E}
\Circle{90.3}{159}{1.1}{dD04}
\Box{94}{155}{44}{\sffamily\fontsize{8.4}{12}\selectfont\color{dD04}D04 模型开放供给}
\definecolor{dD05}{HTML}{537D25}
\Circle{139.3}{159}{1.1}{dD05}
\Box{143}{155}{44}{\sffamily\fontsize{8.4}{12}\selectfont\color{dD05}D05 数据开放供给}
\Rule{22}{166}{188}{line}
\Box{22}{169}{15}{\sffamily\bfseries\fontsize{13}{18}\selectfont\color{gold}X3}
\Box{40}{169}{148}{\sffamily\bfseries\fontsize{12}{18}\selectfont\color{gold}全球连接与网络效应}
\Box{40}{179}{148}{\sffamily\fontsize{10}{15}\selectfont\color{ink}成果如何跨越组织、地域与技术层次形成接续？}
\definecolor{dD06}{HTML}{8E6500}
\Circle{41.3}{194}{1.1}{dD06}
\Box{45}{190}{44}{\sffamily\fontsize{8.4}{12}\selectfont\color{dD06}D06 代码影响力}
\definecolor{dD07}{HTML}{AD4E22}
\Circle{90.3}{194}{1.1}{dD07}
\Box{94}{190}{44}{\sffamily\fontsize{8.4}{12}\selectfont\color{dD07}D07 模型影响力}
\definecolor{dD08}{HTML}{B53863}
\Circle{139.3}{194}{1.1}{dD08}
\Box{143}{190}{44}{\sffamily\fontsize{8.4}{12}\selectfont\color{dD08}D08 数据集影响力}
\Rule{22}{201}{188}{line}
\Box{22}{204}{15}{\sffamily\bfseries\fontsize{13}{18}\selectfont\color{violet}X4}
\Box{40}{204}{148}{\sffamily\bfseries\fontsize{12}{18}\selectfont\color{violet}开源善治与公共价值}
\Box{40}{214}{148}{\sffamily\fontsize{10}{15}\selectfont\color{ink}谁支撑长期维护，开放成果怎样形成共同受益？}
\definecolor{dD09}{HTML}{943B9A}
\Circle{41.3}{229}{1.1}{dD09}
\Box{45}{225}{44}{\sffamily\fontsize{8.4}{12}\selectfont\color{dD09}D09 治理成熟度}
\definecolor{dD10}{HTML}{6146A5}
\Circle{90.3}{229}{1.1}{dD10}
\Box{94}{225}{44}{\sffamily\fontsize{8.4}{12}\selectfont\color{dD10}D10 公共福祉}
\Rule{22}{239}{188}{line}

\Cap{246}{图 1-3　开源发展目标的四支柱与十个观察维度}
\Body{22}{255}{166}{体系能力体现为多类主体与资源之间的持续支撑；公共能力进一步关注更多地区能否掌握、改进和长期使用这些成果。两者需要相互联系、分别观察。}
\Note{22}{272}{166}{来源：2026 版获批研究框架与动态星座 v2.0.0。框架表达发展方向，不等于各维度已完成实测或已实现成效。}
\end{canvas}

% Page 5: public capability conditions, local co-creation and governance.
\begin{canvas}
\Header{1.2 · 发展目标}{teal}\Footer{05}
\Kicker{teal}{1.2 / OPEN RESOURCES, SUSTAINABLE CAPABILITIES}
\Title{让开放成果成为可持续公共能力}
\Subhead{59}{teal}{AI 普惠要连接资源可及、本地共创与长期维护}
\Body{22}{72}{166}{开放使更多人有机会接触和改进技术，但从获得资源到持续使用，仍需许可、成本、人才与组织条件。对于本地语言服务、教育或公共部门应用，关键问题是当地伙伴能否理解系统、参与改进，并在项目交付之后继续维护。}
\Rect{22}{104}{37}{58}{blue!6}
\Box{26}{108}{29}{\sffamily\bfseries\fontsize{19}{23}\selectfont\color{blue}01}
\Box{26}{124}{29}{\sffamily\bfseries\fontsize{12}{17}\selectfont\color{blue}可获得}
\Box{26}{138}{29}{\sffamily\fontsize{9.3}{14}\selectfont\color{ink}资源与权利}
\Box{26}{148}{29}{\sffamily\fontsize{8.1}{12}\selectfont\color{muted}版本、许可、文档}
\LinkLine{60}{132}{64}{132}{line}
\Rect{65}{104}{37}{58}{teal!6}
\Box{69}{108}{29}{\sffamily\bfseries\fontsize{19}{23}\selectfont\color{teal}02}
\Box{69}{124}{29}{\sffamily\bfseries\fontsize{12}{17}\selectfont\color{teal}可复用}
\Box{69}{138}{29}{\sffamily\fontsize{9.3}{14}\selectfont\color{ink}成本与适配}
\Box{69}{148}{29}{\sffamily\fontsize{8.1}{12}\selectfont\color{muted}算力、语言、工具}
\LinkLine{103}{132}{107}{132}{line}
\Rect{108}{104}{37}{58}{gold!6}
\Box{112}{108}{29}{\sffamily\bfseries\fontsize{19}{23}\selectfont\color{gold}03}
\Box{112}{124}{29}{\sffamily\bfseries\fontsize{12}{17}\selectfont\color{gold}可维护}
\Box{112}{138}{29}{\sffamily\fontsize{9.3}{14}\selectfont\color{ink}人员与治理}
\Box{112}{148}{29}{\sffamily\fontsize{8.1}{12}\selectfont\color{muted}责任、资金、安全}
\LinkLine{146}{132}{150}{132}{line}
\Rect{151}{104}{37}{58}{violet!6}
\Box{155}{108}{29}{\sffamily\bfseries\fontsize{19}{23}\selectfont\color{violet}04}
\Box{155}{124}{29}{\sffamily\bfseries\fontsize{12}{17}\selectfont\color{violet}有成效}
\Box{155}{138}{29}{\sffamily\fontsize{9.3}{14}\selectfont\color{ink}场景与结果}
\Box{155}{148}{29}{\sffamily\fontsize{8.1}{12}\selectfont\color{muted}覆盖、质量、公平}
\Box{22}{170}{166}{\sffamily\bfseries\fontsize{8.8}{13}\selectfont\color{navy}图 1-4　开放资源形成公共能力，需要逐项建立条件与结果证据}
\Box{22}{177}{166}{\sffamily\fontsize{7.8}{11}\selectfont\color{muted}前三项关注实施条件，最后一项需由实际场景检验；这是一组分析问题，不是自动生效的因果公式。}

\Label{22}{195}{166}{把共同发展目标转化为可落实的合作条件}
\Body{22}{207}{166}{\textbf{支持本地共创。} 全球南方的语言、文化知识与实际需求应进入问题定义、数据建设和评价过程，使当地伙伴能够参与技术选择与规则形成。}
\Body{22}{226}{166}{\textbf{落实长期组织。} 基金会、企业、研究机构与社区可围绕共同基础设施、维护资金、人才培养和治理责任形成持续合作，具体作用须结合案例检验。}
\Body{22}{245}{166}{\textbf{分别检验条件与结果。} SDGs 提供发展方向，数字公共品标准关注开放与设计条件；成本、服务质量及公平性仍须由真实实施来评价。}
\Note{22}{266}{166}{来源：2026 版第一章条件关系框架；\href{\UNGoals}{联合国可持续发展目标}、\href{\DPG}{DPGA 数字公共品标准}。目标关联与开放条件不构成具体项目的公共成效证明或机构背书。}
\end{canvas}

% Page 6: observation across platforms and object relationships.
\begin{canvas}
\Header{1.3 · 观察方法}{gold}\Footer{06}
\Sec{gold}{1.3　观察方法：跨平台的全球开源全景}{section13}
\Title{将分散的对象，放回协作关系}
\Subhead{59}{gold}{连接人才、组织、项目、代码、模型与数据}
\Body{22}{72}{166}{本报告从 GitHub 的开发协作、HF 的模型与数据资源，以及已具备证据的补充平台出发，联读参与主体、开放资产及其关系。同一项目可能跨越多个仓库，同一模型也可能被不同组织继续加工，单个平台计数难以独自解释完整过程。}
\Box{22}{103}{166}{\sffamily\bfseries\fontsize{10.5}{15}\selectfont\color{teal}以地域与时序，联读主体、资产及其联系}
\Rect{22}{122}{50}{28}{blue!6}
\Box{26}{125}{42}{\sffamily\bfseries\fontsize{11.5}{16}\selectfont\color{blue}人才}
\Box{26}{137}{42}{\sffamily\fontsize{8.5}{13}\selectfont\color{muted}参与、审阅、维护}
\Rect{80}{122}{50}{28}{blue!6}
\Box{84}{125}{42}{\sffamily\bfseries\fontsize{11.5}{16}\selectfont\color{blue}组织}
\Box{84}{137}{42}{\sffamily\fontsize{8.5}{13}\selectfont\color{muted}支持、协调、治理}
\Rect{138}{122}{50}{28}{blue!6}
\Box{142}{125}{42}{\sffamily\bfseries\fontsize{11.5}{16}\selectfont\color{blue}项目}
\Box{142}{137}{42}{\sffamily\fontsize{8.5}{13}\selectfont\color{muted}组织共同工作}
\Rect{22}{177}{50}{28}{teal!6}
\Box{26}{180}{42}{\sffamily\bfseries\fontsize{11.5}{16}\selectfont\color{teal}代码}
\Box{26}{192}{42}{\sffamily\fontsize{8.5}{13}\selectfont\color{muted}开发与工具实现}
\Rect{80}{177}{50}{28}{gold!6}
\Box{84}{180}{42}{\sffamily\bfseries\fontsize{11.5}{16}\selectfont\color{gold}模型}
\Box{84}{192}{42}{\sffamily\fontsize{8.5}{13}\selectfont\color{muted}发布、派生、适配}
\Rect{138}{177}{50}{28}{violet!6}
\Box{142}{180}{42}{\sffamily\bfseries\fontsize{11.5}{16}\selectfont\color{violet}数据}
\Box{142}{192}{42}{\sffamily\fontsize{8.5}{13}\selectfont\color{muted}供给、加工、复用}
\LinkLine{72}{135}{79}{135}{blue}
\LinkLine{130}{135}{137}{135}{blue}
\Box{69}{119}{17}{\sffamily\fontsize{7.1}{10}\selectfont\color{muted}\centering 参与}
\Box{127}{119}{17}{\sffamily\fontsize{7.1}{10}\selectfont\color{muted}\centering 支持}
\draw[teal,line width=.65pt] \Pos{163}{151} -- \Pos{163}{163} -- \Pos{47}{163};
\LinkLine{47}{163}{47}{175}{teal}
\LinkLine{105}{163}{105}{175}{teal}
\LinkLine{163}{163}{163}{175}{teal}
\Box{76}{155}{58}{\sffamily\fontsize{8}{11}\selectfont\color{teal}\centering 形成与维护开放资产}
\draw[gold,line width=.65pt,-{Stealth[length=2mm]}] \Pos{47}{207} -- \Pos{47}{218} -- \Pos{177}{218} -- \Pos{177}{151};
\Rect{72}{214}{89}{8}{white}
\Box{74}{214}{85}{\sffamily\fontsize{8}{12}\selectfont\color{gold}\centering 在其他项目中被复用与再创造}
\Box{22}{229}{166}{\sffamily\bfseries\fontsize{8.8}{13}\selectfont\color{navy}图 1-5　六类对象的关联观察框架}
\Box{22}{236}{166}{\sffamily\fontsize{7.8}{11}\selectfont\color{muted}框架示意：同一实体可对应多个账号、仓库或资源；具体关系须有来源，并登记映射与覆盖状态。}

\Body{22}{249}{166}{跨平台观察的价值，在于识别供给与接续如何连接。身份和关系仅在已核验范围内联读，各类资源保留自己的统计单位；“全域视野”不等于所有对象已被完整测量。}
\Note{22}{271}{166}{来源：2026 版研究框架；\href{\OD}{OpenDigger} 提供数据组织与开放指标基础。平台覆盖、未知地区、身份和项目映射状态随具体图表登记。}
\end{canvas}

% Page 7: accessible explanation, explicit sample/period boundaries.
\begin{canvas}
\Header{1.3 · 观察方法}{gold}\Footer{07}
\Kicker{gold}{1.3 / PARTICIPATION, CONTRIBUTION AND INFLUENCE}
\Title{同一生态，需要不同的观察尺度}
\Subhead{59}{gold}{参与、贡献与影响力，共同解释协作}
\Body{22}{72}{166}{参与说明谁进入共同工作，贡献关注完成了哪些协作，影响力观察这些工作在关系网络中的位置。AI 生成与自动化进一步改变了活动量的含义，因此需要辨别参与主体、协作类型与持续性，不能由提交数量直接判断能力。}
\Rect{22}{105}{57}{92}{paper}
\Box{26}{108}{49}{\sffamily\bfseries\fontsize{9.2}{14}\selectfont\color{navy}项目 P：修复兼容问题}
\Circle{29}{125}{3}{blue!13}
\Box{26}{122.4}{6}{\sffamily\bfseries\fontsize{8}{10}\selectfont\color{blue}\centering 甲}
\Box{36}{122}{39}{\sffamily\fontsize{8.4}{12}\selectfont\color{navy}发现并报告问题}
\LinkLine{29}{129}{29}{135}{blue!40}
\Circle{29}{142}{3}{blue!13}
\Box{26}{139.4}{6}{\sffamily\bfseries\fontsize{8}{10}\selectfont\color{blue}\centering 乙}
\Box{36}{139}{39}{\sffamily\fontsize{8.4}{12}\selectfont\color{navy}提交兼容性修复}
\LinkLine{29}{146}{29}{152}{blue!40}
\Circle{29}{159}{3}{blue!13}
\Box{26}{156.4}{6}{\sffamily\bfseries\fontsize{8}{10}\selectfont\color{blue}\centering 丙}
\Box{36}{156}{39}{\sffamily\fontsize{8.4}{12}\selectfont\color{navy}审阅与验证改动}
\LinkLine{29}{163}{29}{169}{blue!40}
\Circle{29}{176}{3}{blue!13}
\Box{26}{173.4}{6}{\sffamily\bfseries\fontsize{8}{10}\selectfont\color{blue}\centering 丁}
\Box{36}{173}{39}{\sffamily\fontsize{8.4}{12}\selectfont\color{navy}在项目 Q 中复用}
\Box{26}{189}{49}{\sffamily\fontsize{7.2}{10}\selectfont\color{muted}协作示意，不代表实测结果。}
\Box{88}{105}{100}{\sffamily\bfseries\fontsize{11}{16}\selectfont\color{blue}参与}
\Box{88}{114}{100}{\sffamily\fontsize{9.4}{14}\selectfont\color{ink}哪些人持续进入协作？}
\Box{88}{124}{100}{\sffamily\fontsize{8.2}{12}\selectfont\color{muted}人数、留存与角色分布}
\Box{88}{137}{100}{\sffamily\bfseries\fontsize{11}{16}\selectfont\color{teal}贡献}
\Box{88}{146}{100}{\sffamily\fontsize{9.4}{14}\selectfont\color{ink}完成了哪些共同工作？}
\Box{88}{156}{100}{\sffamily\fontsize{8.2}{12}\selectfont\color{muted}开发、讨论、审阅与修复}
\Box{88}{169}{100}{\sffamily\bfseries\fontsize{11}{16}\selectfont\color{gold}影响力}
\Box{88}{178}{100}{\sffamily\fontsize{9.4}{14}\selectfont\color{ink}这些工作怎样被连接和接续？}
\Box{88}{188}{100}{\sffamily\fontsize{8.2}{12}\selectfont\color{muted}协作联系、接续路径与结构位置}
\Box{22}{204}{166}{\sffamily\bfseries\fontsize{8.8}{13}\selectfont\color{navy}图 1-6　同一协作过程，可以从参与、贡献与影响力三个视角观察}
\Box{22}{211}{166}{\sffamily\fontsize{7.8}{11}\selectfont\color{muted}OpenRank 依据定义的协作关系计算网络影响；单次计数、下载或关注均不能直接替代这一观察。}

\Rule{22}{230}{188}{line}
\Label{22}{236}{166}{主统计期：2025-07-01 至 2026-06-30}
\Note{22}{247}{166}{历史比较使用可比窗口；外部平台年度窗和期后观察分别标注。跨平台身份、自动化账号、地区未知与项目边界按实际核验状态处理，未知不按零计算。}
\Note{22}{262}{166}{来源：2026 版第一章与\href{\OD}{OpenDigger／OpenRank 方法材料}。OSDI 将四目标、十方向落实为候选指标；完整定义与可得性在附录 A 及研究记录中维护。下载、关注和派生也须分别解释，不能直接等同公共结果。}
\end{canvas}

% Page 8: four complementary perspectives and research questions, no premature findings.
\begin{canvas}
\Header{1.3 · 全书路径}{gold}\Footer{08}
\Kicker{gold}{1.3 / FROM QUESTIONS TO EVIDENCE AND ACTION}
\Title{从中国出发，理解全球开放协作}
\Subhead{59}{gold}{四类证据共同回答人工智能时代的发展问题}
\Body{22}{72}{166}{中国的模型供给、产业与人才实践，为观察全球开源提供了重要入口。本报告将其放入跨境复用与共同创新中，也关注其他地区的长期积累和全球南方的参与路径。地域特色与时序变化，需要由互补证据共同解释。}
\Rect{22}{100}{166}{19}{blue!5}
\Box{26}{103}{13}{\sffamily\bfseries\fontsize{14}{18}\selectfont\color{blue}02}
\Box{44}{102}{36}{\sffamily\bfseries\fontsize{10.5}{15}\selectfont\color{blue}生态地图}
\Box{44}{111}{36}{\sffamily\fontsize{7.5}{10}\selectfont\color{muted}GOSO}
\Box{85}{105}{99}{\sffamily\fontsize{9.5}{15}\selectfont\color{ink}参与与资源分布在哪里？}
\Rect{22}{122}{166}{19}{gold!5}
\Box{26}{125}{13}{\sffamily\bfseries\fontsize{14}{18}\selectfont\color{gold}03}
\Box{44}{124}{36}{\sffamily\bfseries\fontsize{10.5}{15}\selectfont\color{gold}影响力排行}
\Box{44}{133}{36}{\sffamily\fontsize{7.5}{10}\selectfont\color{muted}GOSR}
\Box{85}{127}{99}{\sffamily\fontsize{9.5}{15}\selectfont\color{ink}哪些主体形成关键影响？}
\Rect{22}{144}{166}{19}{teal!5}
\Box{26}{147}{13}{\sffamily\bfseries\fontsize{14}{18}\selectfont\color{teal}04}
\Box{44}{146}{36}{\sffamily\bfseries\fontsize{10.5}{15}\selectfont\color{teal}技术全景}
\Box{44}{155}{36}{\sffamily\fontsize{7.5}{10}\selectfont\color{muted}GOSL}
\Box{85}{149}{99}{\sffamily\fontsize{9.5}{15}\selectfont\color{ink}技术如何分层、连接与演进？}
\Rect{22}{166}{166}{19}{violet!5}
\Box{26}{169}{13}{\sffamily\bfseries\fontsize{14}{18}\selectfont\color{violet}05}
\Box{44}{168}{36}{\sffamily\bfseries\fontsize{10.5}{15}\selectfont\color{violet}模型生态}
\Box{44}{177}{36}{\sffamily\fontsize{7.5}{10}\selectfont\color{muted}GOSM}
\Box{85}{171}{99}{\sffamily\fontsize{9.5}{15}\selectfont\color{ink}模型如何供给、复用与开放？}
\Box{22}{194}{166}{\sffamily\bfseries\fontsize{8.8}{13}\selectfont\color{navy}图 1-7　第一章提出问题与方法，第二至第五章提供四类互补证据}

\Label{22}{207}{166}{贯穿全书的四组研究问题}
\Label{22}{219}{14}{\textcolor{blue}{01}}
\Body{37}{218}{151}{全球协作怎样形成跨层支撑与持续接续的体系能力？}
\Label{22}{231}{14}{\textcolor{teal}{02}}
\Body{37}{230}{151}{中国模型供给如何连接工具链、制度建设与全球社区？}
\Label{22}{243}{14}{\textcolor{gold}{03}}
\Body{37}{242}{151}{全球南方如何通过本地共创和组织支持形成公共能力？}
\Label{22}{255}{14}{\textcolor{violet}{04}}
\Body{37}{254}{151}{中国不同产业、人才群体和城市怎样形成各自优势？}
\Note{22}{270}{166}{第六章依据前文经验证的发现，提炼关键洞察、分层时间前瞻与行动建议；各研究视角互补，不强制一一对应。}
\end{canvas}
\end{document}
